% ------------------------
% Resume - Machine Learning Engineer
% Author : Dipesh Maindolia
% ------------------------

\documentclass[letterpaper,8pt]{article}

\usepackage[empty]{fullpage}
\usepackage[margin=0.24in, top=0.1in, bottom=0.1in]{geometry}
\usepackage{titlesec}
\usepackage{enumitem}
\usepackage[hidelinks]{hyperref}
\usepackage{xcolor}
\usepackage{fontenc}

\pagestyle{empty}
\sloppy
\setlength{\parfillskip}{0pt plus 1fil}

\urlstyle{same}
\linespread{0.80}
\setlength{\parindent}{0pt}

\setlength{\tabcolsep}{0in}
\setlength{\parskip}{0pt}

\definecolor{linkblue}{RGB}{26, 86, 160}

\setlist[itemize]{
  leftmargin=0.18in,
  noitemsep,
  topsep=1pt,
  parsep=0pt,
  partopsep=0pt
}

% ---------- SECTION STYLE ----------
\titleformat{\section}{
  \vspace{0pt}\scshape\raggedright\normalsize\bfseries
}{}{0em}{}[\titlerule \vspace{0pt}]

% ---------- CUSTOM COMMANDS ----------
\newcommand{\resumeSubheading}[4]{
  \begin{tabular*}{\textwidth}{l @{\extracolsep{\fill}} r}
    \textbf{#1} & \small\textit{#2} \\
    \textit{\small#3} & \textit{\small #4} \\
  \end{tabular*}
}

\newcommand{\resumeProject}[3]{
  \begin{tabular*}{\textwidth}{l @{\extracolsep{\fill}} r}
    \textbf{#1} \small{-- #2} & \small\href{#3}{\color{linkblue}\underline{GitHub}} \\
  \end{tabular*}
}

% ---------- DOCUMENT ----------
\begin{document}

% ---------- HEADER ----------
\begin{center}
    {\LARGE \textbf{Dipesh Maindolia}} \\[2pt]
    {\small \textit{Machine Learning Engineer}} \\[2pt]
    {\small
    Ramnagar, Uttarakhand, India \textbar\
    +91 7017982390 \textbar\
    \href{mailto:dipeshmaindolia000@gmail.com}{\color{linkblue}dipeshmaindolia000@gmail.com}
    \textbar\
    \href{https://linkedin.com/in/dipeshmaindolia}{\color{linkblue}LinkedIn}
    \textbar\
    \href{https://github.com/dipeshmaindolia001}{\color{linkblue}GitHub}
    }
\end{center}

\vspace{-15pt}

% ---------- SUMMARY ----------
\section{Professional Summary}
\small{Machine Learning Engineer with hands-on experience building end-to-end ML and GenAI pipelines --
from data ingestion and feature engineering to model deployment. Skilled in Python, XGBoost/scikit-learn,
LLM-based application development, and cloud data pipelines on AWS, with working knowledge of Docker and
MCP servers for building and deploying model-serving applications. GATE CSE and GATE DA qualified.}

\vspace{-15pt}

% ---------- SKILLS ----------
\section{Technical Skills}
\begin{itemize}
    \item \textbf{Programming \& Query Languages:} Python, SQL, C++
    \item \textbf{Machine Learning:} scikit-learn, XGBoost, Feature Engineering, Hyperparameter Tuning (Bayesian Optimization), Model Evaluation
    \item \textbf{Generative AI:} LLM Application Development, Prompt Engineering, Few-Shot Learning, LangChain, Groq API (Llama 3.2), MCP Servers
    \item \textbf{MLOps \& Deployment:} Docker, Streamlit, AWS (S3, Lambda, Glue, Step Functions, SageMaker), Model Serving, CI-Ready Pipelines
    \item \textbf{Data Engineering:} ETL / Data Pipelines, Medallion Architecture, Data Cleaning, Data Validation, Pipeline Monitoring
    \item \textbf{Core Competencies:} Statistical Analysis, Data Visualization, Problem Solving, Documentation
\end{itemize}

\vspace{-15pt}

% ---------- EXPERIENCE ----------
\section{Experience}

\resumeSubheading{Python \& Data Science Engineer}{August 2026 -- Present}{Blazop}{}
\begin{itemize}
    \item Configured and integrated MCP (Model Context Protocol) servers for Meraki, Symphony AI, and other enterprise platforms, exposing their APIs as structured tools so AI agents can query and act on them reliably.
    \item Built a type-agnostic AWS Load Balancer (ELBv2) adapter for the LBaaS platform supporting ALB, NLB, and GWLB, with create, delete, modify-subnet, and status operations driven by a manifest passed directly to Boto3.
    \item Added multi-tenant support to the LBaaS adapter through cross-account STS AssumeRole, load balancer ARN resolution by name, optional S3 access-log configuration, and JSON-safe responses with structured logging and error handling.
    \item Automated Terraform provisioning via GitHub Actions: built a GitHub access layer that validates order data against JSON schemas, commits it to auto-created branches, and triggers workflow\_dispatch runs to provision infrastructure.
    \item Calculated token consumption for Claude and OpenAI models and built token management to track usage, support GPU capacity planning, and control costs.
\end{itemize}

\vspace{1pt}

\resumeSubheading{Product Analyst Intern}{April 2026 -- June 2026}{Interviewkit.ai (Remote)}{}
\begin{itemize}
    \item Collaborated directly with the founder on product strategy, feature planning, and execution, translating usage data and user feedback into prioritized improvements to platform workflows and engagement systems.
    \item Built a REST API editor and an intelligent job recommendation module using keyword extraction, structured filtering, and profile-matching techniques to surface the most relevant listings for each user.
    \item Developed a resume-based system to auto-generate personalized technical interview questions and detailed, role-specific model answers for candidates preparing across multiple domains.
\end{itemize}

\vspace{-15pt}

% ---------- PROJECTS ----------
\section{Projects}

\resumeProject{Flight Price Prediction using AWS SageMaker}{Python, XGBoost, scikit-learn, Streamlit}{https://github.com/dipeshmaindolia001/sagemaker-flight-prices-prediction}
\begin{itemize}
    \item \textbf{End-to-End ML Pipeline:} Built a complete flight price prediction system on AWS SageMaker, covering preprocessing, feature engineering, and hyperparameter tuning using XGBoost with Bayesian optimization.
    \item \textbf{Feature Engineering:} Engineered features from raw flight data (airline, route, timestamps, duration) using scikit-learn Pipelines, FeatureUnion, and custom transformers.
    \item \textbf{Deployment:} Deployed a Streamlit web application to Streamlit Cloud, enabling users to input flight details and receive real-time price predictions.
\end{itemize}

\vspace{1pt}

\resumeProject{LinkedIn Post Generator -- Gen AI (Llama 3.2)}{Python, LangChain, Groq API, Streamlit}{https://github.com/dipeshmaindolia001/LinkedIn-Post-Generator-Gen-AI-Using-Llama3.2}
\begin{itemize}
    \item \textbf{GenAI Pipeline:} Built a two-stage Generative AI pipeline that extracts topic, language, and length metadata from a user's past LinkedIn posts, then uses Llama 3.2 (via the Groq API) to generate new posts that match their personal writing style.
    \item \textbf{Prompt \& Retrieval Design:} Implemented few-shot prompting by retrieving stylistically similar past posts for the selected topic, language, and length, and used them to guide the LLM's tone, structure, and output quality.
    \item \textbf{Application Layer:} Deployed an interactive Streamlit application allowing users to select topic, length, and language, and generate on-brand posts on demand with minimal manual editing.
\end{itemize}

\vspace{1pt}

\resumeProject{YouTube Trending Data Pipeline on AWS}{AWS S3, Lambda, Glue, Step Functions}{https://github.com/dipeshmaindolia001/YouTube-Trending-Data-Pipeline-on-AWS}
\begin{itemize}
    \item \textbf{Pipeline Design:} Built a fully automated end-to-end data pipeline using AWS S3, Lambda, and Glue with the YouTube Data API to ingest, clean, transform, and store data.
    \item \textbf{Architecture \& Orchestration:} Implemented Medallion Architecture (Bronze/Silver/Gold) with AWS Step Functions for reliable orchestration, scheduling, and retry logic.
    \item \textbf{Monitoring \& Reliability:} Configured automated alerts across pipeline stages to ensure consistent, analysis-ready data with minimal manual intervention.
\end{itemize}

\vspace{-15pt}

% ---------- EDUCATION ----------
\section{Education}

\begin{tabular*}{\textwidth}{l @{\extracolsep{\fill}} c r}
\textbf{B.Tech in Information Technology} & \textbf{CGPA: 7.56} & \textbf{2022 -- 2026} \\
\textit{College of Technology, Pantnagar} & & \\[1pt]

\textbf{Senior Secondary (Class XII)} & \textbf{92\%} & \textbf{2021} \\
\textit{St. Joseph Convent School} & & \\[1pt]

\textbf{High School (Class X)} & \textbf{92\%} & \textbf{2019} \\
\textit{USR Indu Inter College} & & \\
\end{tabular*}

\vspace{-16pt}

% ---------- POSITION OF RESPONSIBILITY ----------
\section{Position of Responsibility}

\textbf{Convocation Digital Card Team Lead}
\begin{itemize}
    \item Led the end-to-end digital invitation and QR-based security verification system for the university convocation, coordinating a cross-functional team and eliminating paper-based workflows to contribute Rs.~1.5 lakh in cost savings via real-time entry validation and audit logs.
\end{itemize}

\vspace{-16pt}

% ---------- ACHIEVEMENTS ----------
\section{Achievements}

\begin{itemize}
    \item \textbf{GATE:} Cleared \textbf{GATE CSE 2026} (94.3 percentile) and \textbf{GATE DA 2025} (86 percentile), qualifying among the top rankers nationally.
    \item \textbf{Certification:} Completed the \textbf{Foundations of Data Science and Applications} certification from the \textbf{Indian Institute of Technology (IIT) Madras}.
    \item \textbf{EdTech Initiative:} Built and scaled
    \href{https://pantnagar-pathshala.vercel.app/}{\color{linkblue}\underline{Pantnagar Pathshala}} with an associated
    \href{https://www.youtube.com/@dipeshmaindolia}{\color{linkblue}\underline{YouTube channel}} that crossed
    \textbf{1,000+ subscribers}.
\end{itemize}

\end{document}
